\chapter{Lifting in characteristic two: a quadratic label budget}\label{tx:sec:quadratic-route}
\label{tx:ch:quadratic}
\section*{The property we need}
We remain in characteristic two. The block has $q=2^a$ labels,
and $z=\sum_i c_i(\sigma_i-I)$ is its square-zero translation
operator. A common witness is a group-algebra element $w$ satisfying
$zwz=z$ on every injective-tuple space $P_b$ with $b\le r$.
The block stays fixed throughout the following induction; the
retained labels are the resource we count.

We now construct the witness used in the previous chapter. The first
construction is intentionally generous with labels: at tuple length
$b$, it allocates $b$ labels which have not been used before. It
therefore proves exact deletion with a quadratic retained-label
budget. This gives a complete route to a square-root degree lower
bound and isolates exactly where that route loses a factor of $r$.

Keep the progression picture in view. Given a degree-$d$ proof at
one size eligible for this construction, we obtain a degree-$d$ proof
at every later size separated by the same block $q$. Thus $d$ is constant along
that progression. The quadratic cost below governs which starting
proofs are eligible; its inverse is the square-root degree lower bound.

The induction maintains one group-algebra element, not a list of
independently chosen inverses. At the start of stage $b$, this element
satisfies $zwz=z$ on every shorter tuple module. We must repair length
$b$ while preserving all those previous identities. Three steps do
that: normalise the old element, localise its residual to tuples
without the fresh labels, and simulate a position-compatible inverse
by permutations involving those fresh labels.

Keep the order of those steps. A correct partial inverse of $z$ on one
vector space does not automatically give a partial inverse for the old residual, and an
operator which works after projecting back to old labels need not
work on all larger tuples. Normalisation settles the first issue;
localisation settles the second.

\section{Fresh-label construction}\label{tx:sec:quadratic}
We now prove that a support-$r$ witness exists in characteristic two when
\[n\ge T(r):=1+2+\cdots+r=\frac{r(r+1)}2,\qquad r<q.\]
This proof is independent of the rook theorem used later. Its ingredients are finite algebra, a position-compatible auxiliary inverse, and an explicit way to simulate that inverse using fresh labels.

The phrase ``generic coefficients'' below refers to the independent
block coefficients from Chapter~\ref{tx:ch:square-zero}. We choose
them in a finite field so that the subsequent ring calculation takes
place in a finite set. If $\Omega$ is the current finite label set,
$F[S_\Omega]$ has only $|F|^{|S_\Omega|}$ elements: each permutation
has one coefficient in $F$. The finite field is a construction device,
not a restriction on the field of the original polynomial request.

Choose the generic coefficients in a finite field $F$ of characteristic two, large enough to contain $a$ independent elements. We may work inside the algebraic closure of the original field; the degree question descends by the field-extension argument. The finiteness of $F$ will be used only to normalise an element of a finite group algebra.

\subsection{Normalisation of the residual}
Suppose $w_0$ works on shorter tuples but not yet on length $b$. Put $y=z-zw_0z$. An auxiliary linear map $h$ may satisfy $zhz=z$ on the larger tuple space. It is tempting to conclude $yhy=y$. In general the expansion contains an extra term $zw_0zw_0z$, and this need not collapse. Before correcting the error we arrange the additional identity $w_0zw_0=w_0$ in the actual group algebra. This normalisation was a genuine mathematical step in the earlier construction.

Here \emph{normalisation} is an algebraic replacement: it modifies
a ring element to satisfy the
additional algebraic identity needed for the residual calculation.
A unital ring has addition, associative multiplication and an identity
$1$; its multiplication need not commute. A ring representation applies
its elements as operators, respecting addition, composition and the
identity. Thus an identity proved in the ring holds in all its
representations at once. In the lemma, $v$ is a ring element and
the letters $i,\ell,h$ in the proof are integers recording powers;
this local $h$ is not the auxiliary linear operator used elsewhere.

\begin{lemma}[Normalisation in a finite ring]\label{tx:lem:normal}
Let $v,z$ belong to a finite unital ring. There is an element
$v^\natural$ in the same ring with
$v^\natural z v^\natural=v^\natural$, such that every ring
representation in which $zvz=z$ also satisfies
$zv^\natural z=z$.
\end{lemma}
\begin{proof}
Put $x=vz$. Finiteness gives $x^{i+\ell}=x^i$ for some $i\ge0$
and $\ell>0$. Multiplication by powers of $x$ gives
$x^{j+\ell}=x^j$ whenever $j\ge i$. Choose $h\ge i$ such that
$h+1$ is divisible by $\ell$. Then $2h+1-h=h+1$ is a multiple
of $\ell$, so $x^{2h+1}=x^h$. Define $v^\natural=x^hv$.
Associativity, without commutation of factors, gives
\[
 v^\natural z v^\natural
      =x^h(vz)x^hv=x^{2h+1}v=x^hv=v^\natural.
\]
In a representation with $zvz=z$, induction gives
$z(vz)^j=z$ for every $j\ge0$. Hence
\[
                     zv^\natural z=z(vz)^{h+1}=z.
\]
The exponent was selected inside the original finite ring. The
normalisation therefore preserves the witness relation simultaneously
in every representation where it was already true.
\end{proof}
In the construction below we denote the normalised old witness by $b_0$. It is one ring element shared by every shorter tuple representation.
The lemma preserves $zvz=z$ wherever it was already valid. It does
not assert that a formerly unsuccessful tuple length has been repaired.
The new identity $b_0zb_0=b_0$, obtained in the ring itself, is what
makes the following sandwich of the residual work.

\begin{lemma}[Normalised residual]\label{tx:lem:residual}
If $b_0zb_0=b_0$ and $zhz=z$, then, for $y=z-zb_0z$,
\[yhy=y.\]
\end{lemma}
\begin{proof}
Expand without commuting any factors:
\begin{align*}
yhy&=zhz-zhzb_0z-zb_0zhz+zb_0zhzb_0z\\
&=z-2zb_0z+zb_0zb_0z\\
&=z-zb_0z=y.
\end{align*}
The last equality uses the normalisation. This identity works in any ring; characteristic two is not needed for this particular calculation.
\end{proof}

\subsection{A residual calculation over \texorpdfstring{$\mathbb F_4$}{F4}}
The notation $E_{ij}$ denotes a matrix unit: it sends basis vector
$e_j$ to $e_i$ and sends the other basis vectors to zero. Thus
$E_{12}(s,t)=(t,0)$ and $E_{21}(s,t)=(0,s)$ on a two-dimensional
coordinate space. Applying $E_{12},E_{21},E_{12}$ successively
returns $E_{12}(s,t)$, which makes the partial-inverse identity
in the calculation below transparent.

The normalisation is not optional algebraic tidying. Let
$\omega^2+\omega+1=0$, let $z=E_{12}$ on $\mathbb F_4^2$, and
let $h=E_{21}$. Then $zhz=z$. If we take the unnormalised old
operator $w_0=\omega E_{21}$, its residual is
\[
                   y=z-zw_0z=(1-\omega)z=\omega^2z.
\]
But $yhy=\omega^4z=\omega z\ne\omega^2z=y$. Thus a partial inverse
for $z$ does not automatically serve as an inverse for $y$. The
finite-ring lemma repairs exactly the missing identity used in the
expansion of $yhy$.

This example is a two-dimensional operator calculation explaining
one algebraic step. It is not being used as a label-permutation
construction or as an asymptotic counterexample. The group-algebra
construction obtains its normalisation in the finite ring
$F[S_\Omega]$ itself, which is why the corrected identity holds in
all the required modules at once.

\subsection{An auxiliary inverse compatible with tuple positions}
Let $\Omega$ be the current label set, consisting of the block $B$ and the retained labels already allocated. On $Q=P_b(\Omega)$, partition the tuples according to the $A$-orbit of their \emph{underlying label set}. A set meeting $B$ has proper stabiliser because it has size $b<q$. Choose the functional $\ell$ and coordinate $\pi$ as in Section~\ref{tx:sec:block2}, now on these sets. Let $D$ read $\pi$ from the underlying set, and use the same $W$ and nonzero $\alpha$ as before. The choices and the scalar $\alpha$ are made separately on each orbit; the resulting operators are combined on their direct sum. On each moved orbit,
\[Dz-zD=W,\quad Wz=zW,\quad W^2=\alpha^2I,\quad z^2=0.\]
Multiplying the commutator by $z$ shows $zDz=zW$. Thus
\[h=\alpha^{-2}WD\quad\text{satisfies}\quad zhz=z.\]
On tuples wholly outside $B$, put $h=0$, since $z=0$ there. Because $D$ depends on the underlying set and translations act on labels, $h$ commutes with every permutation of tuple positions. We have not claimed that $h$ is a group-algebra element. It is just an explicitly constructed linear operator.
Here and in this section, ``inverse'' is shorthand for the partial
identity $zhz=z$, not for $zh=hz=I$.

\subsection{Fresh labels simulate any position-compatible operator}
Write $C=\{\kappa_1,\ldots,\kappa_b\}$ for the fresh labels.
The old tuple space sits inside the larger one by including its basis
vectors; projection back discards the basis vectors using a fresh label.
The composite of that inclusion and projection on the old space is
the identity. We want only the old-to-old part of a larger permutation
combination, which is the meaning of the compressed operator below.

Adjoin $b$ new labels $\kappa_1,\ldots,\kappa_b$ to $\Omega$. Let $i:P_b(\Omega)\to P_b(\Omega\cup C)$ include the old tuples and let $p$ discard every tuple using a fresh label. We prove that any operator $h$ commuting with position permutations has the form
\begin{equation}\label{tx:eq:compression}h=pci\end{equation}
for some $c\in F[S_{\Omega\cup C}]$.
The order of operations in $pci$ is: include the old input, apply
$c$, then project the result to old outputs. The symbol $p$ here
names this projection, not the prime characteristic. The formula
does not say that $c$ preserves the old space or agrees with $h$
on every larger tuple; it identifies its old-to-old matrix corner.

Fix a $b$-element set $U=\{u_1,\ldots,u_b\}\subseteq\Omega$, with one chosen ordering. Form
\[E_U=\prod_{j=1}^b\bigl(1-(u_j\ \kappa_j)\bigr).\]
The transpositions are disjoint. On an old tuple missing some $u_j$, the corresponding factor kills the tuple, and the other factors cannot introduce $u_j$. On a tuple using all the $u_j$, the identity term leaves the tuple old, while every nonidentity term introduces a fresh label. Therefore $pE_Ui$ is precisely the projector onto the orderings of $U$.

The $b!$ orderings of $U$ form a \emph{free orbit} for position
permutations: exactly one reordering carries a chosen ordering to
any other. Consequently a position-compatible map is determined
on that orbit by its value at one ordering. All other values are
obtained by the same position reordering of the output. This is
why one can prescribe the map without averaging or dividing by $b!$.

For an ordered output tuple $v$, choose a permutation $g\in S_\Omega$ taking $(u_1,\ldots,u_b)$ to $v$. Extend $g$ by fixing fresh labels. Then $pgE_Ui$ takes that input ordering to $v$, takes its position permutations to the corresponding permutations of $v$, and kills all other underlying sets. The orderings of each $U$ form a free orbit for $S_b$. Consequently an operator commuting with $S_b$ is completely determined by its value on one ordering for each $U$. Expanding those values as linear combinations of output tuples and using the operators $pgE_Ui$ proves \eqref{tx:eq:compression}. There is no division by $b!$.

\begin{example}[Two registers]
For $U=\{u_1,u_2\}$,
\[E_U=1-(u_1\kappa_1)-(u_2\kappa_2)+(u_1\kappa_1)(u_2\kappa_2).
\]
On $(u_1,u_2)$, projection $p$ kills the last three terms. On $(u_1,v)$ with $v\ne u_2$, the identity contribution is cancelled by the second swap factor, which fixes that tuple. The detector recognises the set $U$, while leaving the two possible orderings available for a position-compatible output rule.
\end{example}

\subsection{Why compression is enough only after localisation}
Equation~\eqref{tx:eq:compression} asserts $h=pci$ on the old tuple
space. It does not assert that $c$ acts like $h$ on every tuple in
the larger universe. Terms of $c$ may introduce fresh labels. This
would normally be a serious problem: a proof must preserve all basis
objects, not just an arbitrarily selected corner of a matrix.

The residual has a special support in the larger tuple space. It is
zero on every summand with a fresh label, because such a summand uses
fewer than $b$ old labels and the old witness is already valid at
those shorter lengths. Therefore the two residual factors in
$ycy$ project to the old corner before and after the middle operator.
The identity $h=pci$ is exactly the information the sandwich needs.
The next calculation verifies this point as an equality of maps,
rather than assuming that compression itself is a global realisation.

In matrix terms, order the basis with old tuples first and tuples
using fresh labels second. Under the induction hypothesis the
calculation below gives
\[
                       y=\begin{pmatrix}y_Q&0\\0&0\end{pmatrix}.
\]
In $ycy$ the rightmost $y$ removes fresh input directions, and the
leftmost $y$ removes fresh output directions. Only the old-to-old
corner $pci=h$ of the middle operator matters. The zeros here come
from validity of the old witness at shorter lengths, not from a
claim that an arbitrary permutation combination has this block form.

\subsection{The induction, including localisation of the error}
For support zero take $w_0=0$: $z$ is zero on constants. Suppose a witness exists on the block and $T(b-1)$ retained labels, collectively $\Omega$. Normalise it in the finite ring $F[S_\Omega]$ to $b_0$. It still works on all tuple lengths below $b$ and satisfies $b_0zb_0=b_0$ in the ring itself. Set $y=z-zb_0z$.

Adjoin $b$ fresh labels. Decompose $P_b(\Omega\cup C)$ according to the positions occupied by fresh labels and their actual values. The group $S_\Omega$ fixes each fresh label, so every summand with at least one fresh label is a copy of $P_j(\Omega)$ for $j<b$. On it $y=0$ by the old witness relation. The summand with no fresh labels is $Q=P_b(\Omega)$. Thus, as actual maps,
\[y=i y_Q p.\]
This localisation is what allows a compressed inverse to suffice. Construct the auxiliary $h$ on $Q$ above. Lemma~\ref{tx:lem:residual} gives $y_Qhy_Q=y_Q$. Choose $c$ with $pci=h$. Then
\[ycy=i y_Q(pci)y_Qp=i y_Qp=y.\]
Define the new group-algebra element
\[w_b=b_0+(1-b_0z)c(1-zb_0).\]
Since $z(1-b_0z)=(1-zb_0)z=y$,
\[zw_bz=zb_0z+ycy=zb_0z+y=z\]
on length $b$. On every shorter tuple space the same fresh-label decomposition makes $y=0$, so the relation remains valid there too. The new cost is exactly $b$ retained labels, giving
\[T(b)=T(b-1)+b.\]
If still more retained labels are present, extend the constructed permutations by fixing them. Decomposing according to those outside labels reduces every tuple of length at most $b$ to a shorter tuple on the constructed labels. Thus the same witness works for all $n\ge T(b)$.

\begin{theorem}[Quadratic retained-label budget]
For every $r<q=2^a$ and $n\ge r(r+1)/2$, there is a common witness $w$ with $zwz=z$ on all tuple modules of length at most $r$. Therefore deletion is exact on the support-$r$ class. It transports the
entire degree-$d$ PC consequence space with $r=kd+s$, beyond the fixed
template threshold. For prescribed indexed targets use the budget
$r=2kd+s+t$ to include the target map and the doubled-support detector.
\end{theorem}
Every construction needed for this theorem has now been supplied. In particular, the proof is not an inference from witnesses found for small tuple lengths.


\begin{remark}[The square-root degree bound]
For either prescribed-target budget, the construction gives an onset
$O(d^2)$ for degree-$d$ transfer.
Inverting that onset gives an $\Omega(\sqrt n)$ lower bound when no
finite degree recurs, and on the hard residues of the periodic case.
The loss is precisely the cumulative allocation $T(r)=r(r+1)/2$.
The following construction replaces this sum by the maximum $2r-1$,
using one retained pool at every layer.
\end{remark}

\section{The first three induction stages}
At length zero there is one basis vector, the empty tuple. Every
permutation fixes it, so $z=0$ and $w_0=0$ is a valid witness. No
retained label is needed.

At length one, the auxiliary operator recognises the moved label
orbits. One fresh retained label realises its compressed inverse.
The witness can therefore be built on the block plus one retained
label. For a block of two labels, the previous chapter exhibited
this witness explicitly as a transposition between one block label
and the retained label.

At length two, normalise that same witness first. Allocate two
additional retained labels $\kappa_1,\kappa_2$. A tuple containing
either fresh label has at most one old entry; the residual is zero
on it. Only the tuples with both entries old require repair. The
detector
\[
                 (1-(u_1\ \kappa_1))(1-(u_2\ \kappa_2))
\]
selects an old two-element set. Combining its translates realises
any position-compatible inverse on the old corner. The total number
of retained labels is now $1+2=3$.

At length three, the same argument spends three further labels.
Tuples with one, two or three of those fresh labels reduce to old
tuple lengths two, one or zero, respectively. All three pieces
already have zero residual. The only unsolved piece consists of
old three-tuples, and a detector with three disjoint swaps compresses
its auxiliary inverse. The retained total is $1+2+3=6$.

The block size must exceed the largest tuple length being treated.
For example, one may choose a block of size four when explaining
stages one through three. A two-label block illustrates stage one,
but cannot be reused as a valid block for stage three. Retained
labels and block labels have different jobs in this count.

\begin{center}
\begin{tabular}{@{}c|rrrr@{}}
Tuple length $r$ &1&2&3&4\\
Fresh-label total $T(r)$ &1&3&6&10\\
Linear total proved later $2r-1$ &1&3&5&7
\end{tabular}
\end{center}
The first strict improvement is at length three. The linear method
will not allocate the third batch separately: it will exploit the
residual's zero marginals and reuse a common retained pool.

\section{From quadratic label cost to square-root degree hardness}
Return now from tuple spaces to one fixed polynomial request and one
proof system, NS or PC. Its profile $h(n)$ is the least degree for
one fixed target orbit, or the maximum of the least degrees for an
all-target request. It is infinity if any requested target is
underivable; an empty request has value zero. A finite value recurs if that exact value is
attained at infinitely many universe sizes. The block construction
will force recurrence whenever such a value occurs beyond its
eligible starting universe size.

Let $r(d)=\alpha d+\beta$, where
$(\alpha,\beta)=(k,s)$ for an invariant target or
$(\alpha,\beta)=(2k,s+t)$ for an indexed request. The fixed starting
universe size $n_*$ includes degree stabilisation of the targets. A convenient
safe onset for this preliminary construction is
\[
 N_{\rm quad}(d)=
  \max\{n_*,\,2r(d)+1,\,T(r(d))\}\qquad(d\ge1).
\]
The extra linear term is harmless and makes this onset comparable
with the eventual linear statement. Exactness of the quadratic
construction, together with the target-map and detection arguments,
gives degree-threshold periodicity beyond this onset. Take the least
power of two $Q(d)>r(d)$; then $Q(d)\le2r(d)$.

Here $\alpha,\beta$ are integer coefficients of the support budget,
unrelated to the nonzero field scalar $\alpha$ in the orbit
calculation. The onset estimates concern positive degree budgets.
Degree zero uses its previously established fixed predicate beyond
$n_*$; it therefore cannot be a nonrecurring value on a hard tail.
The hard-value inversion below concerns positive finite values,
and the unbounded recurring values chosen for peaks are positive too.

Suppose a profile has no recurring finite degree. If $h(n)=e<\infty$
and $n\ge N_{\rm quad}(e)$, threshold periodicity at all budgets
through $e$ repeats that exact value on an infinite arithmetic
progression. This is a contradiction. Thus every finite value must
satisfy $n<N_{\rm quad}(e)$. Since
$N_{\rm quad}(e)\le A(e+1)^2$ for a constant $A$ depending only
on the fixed request, it follows that
\[
                          h(n)>\sqrt{n/A}-1
\]
whenever $n$ is beyond the fixed initial range. The assertion also
holds when $h(n)=\infty$. The same argument works on hard residues
in a fixed periodic easy/hard split: a value attained beyond its
onset would recur and would therefore be one of the bounded easy
values.

For a fluctuating profile, choose recurring values $e$ tending to
infinity. A representative of its recurring residue can be selected
in
\[
       N_{\rm quad}(e)\le n<N_{\rm quad}(e)+Q(e).
\]
Both the onset and the period are bounded by a constant times
$(e+1)^2$, so these representatives are finite degree peaks with a square-root lower bound, $e\ge c\sqrt n-O(1)$, for a fixed $c>0$. They are not an
all-size lower bound; low recurring values remain at arbitrarily
large sizes.

Nothing in this argument makes the degree of an individual proof
shrink. At fixed $d$ we compare derivability at two universe sizes.
The lower bound follows by comparing that \emph{fixed budget} with
its allowed universe-size onset. The preliminary route just proved
is in characteristic two. The later prime-characteristic chapter
will give the full linear result directly over every fixed positive
characteristic.
